% !TeX root = ../papers/nc-p3.tex
\section{五次元交叉交換子の二つの逆像と中心一次枝への分岐}
\label{sec:marked-five-dimensional-bracket-fiber}
\begin{theorem}\label{thm:marked-five-dimensional-bracket-fiber}
同時に二つの座標枠を取って次の形にできる非零交叉交換子テンソルを考える。
$C_{ii}$ は跡零対角空間を張り、$i\ne j$ では $C_{ij}$ は残る添字 $k$ の
$K_k=E_{k+1,k+2}-E_{k+2,k+1}$ の非零倍である（添字は三巡回）。
この形の実現可能な $[C]$ は、独立な二つの共役まで
\[
 C(t)_{ij}=[K_i+tD_i,K_j-tD_j],\qquad
 D_i=E_{i+1,i+2}+E_{i+2,i+1},\qquad t(t^2-1)\ne0
\]
で尽きる。$[C(t')]$ と $[C(t)]$ の軌道が一致する必要十分条件は
$t'^2\in\{t^2,t^{-2}\}$ である。

この軌道上の二切断窓を\emph{全て}復元すると、スカラーと印を忘れた正則な添字保持同型類は以下で尽くされる。
$x=t^2$ とし、$W(t)$ を命題\ref{prop:split-mixed-transverse-curve}の既知の窓とする。
\begin{enumerate}
\item $x\ne-1$ のとき、二つの最小周期1の類 $A(t),A(-1/t)$ が常にある。
$x\ne-3,-1/3$ なら、これらと相異なる最小周期2の二類 $W(t),W(-1/t)$ もある。
$x=-3,-1/3$ では後二候補だけが非正則になる。
\item $x=-1$ なら全候補が一つの正則な周期1の類 $A(t)=W(t)$ になる。
\item 直線の端 $x=1$ で非対角の全 $C_{ij}$ が消えた指定値も含めると、
そこにある正則な添字保持同型類は $W(1)$ ただ一つである。
\end{enumerate}
ここで $A(t)$ は以下の三生成量子多項式環の中心一次拡大である。
この枝全体で $A(t)\simeq A(s)$ と $s=\pm t$ は同値である。
$x=0$ にある交代三空間は五次元像を持たず本定理の軌道条件の外にある。
\end{theorem}
\begin{proof}[軌道の認識と二枚の逆像]
$C_{ij}=c_{ij}K_k$ と書く。$C_{ij}=[X_i,Y_j]$ なら前節の跡の線型条件が成り立つ。
$\operatorname{tr}(X_i C_{ij})=0$ により $X_i$ の交代射影は $a_iK_i$ となる。
また $\operatorname{tr}(X_i C_{jj})=-\operatorname{tr}(X_j C_{ij})=0$ である。
右辺の消失は $i\ne j$ なら $C_{ij}$ が $C_{ji}$ に比例することによる。
従って全 $X_i$ の対角は零である。同様に $Y_i$ の交代射影は $b_iK_i$、対角は零となる。
相異なる $i,j,k$ の跡条件は $a_i c_{kj}+a_k c_{ij}=0$ である。
一つの $a_i$ が零なら全て零となる。全て零の場合、$Y$ の交代部分も零なら
交叉交換子の対角成分が出ない。$b_i$ の一つが非零なら同じ跡条件で全て非零となり、
対称成分を $X_i=\sum h_{li}D_l$ と書いて $[X_i,Y_j]$ の対角成分を $i\ne j$ で取ると
$h_{ji}=0$。残る $X_i=h_{ii}D_i$ に対し非対角の対称成分が
$h_{ii}b_j[D_i,K_j]=0$ となり、やはり解がない。$b_i$ についても同様である。
従って二つの三つ組とも全て非零で、上の三式の積から
$c_{01}c_{12}c_{20}=-c_{10}c_{21}c_{02}$ である。
入力側の対角共役の比の平方を選べば $C_{ji}=-C_{ij}$ を揃えられる。

$C_{ij}=\epsilon_{ijk}\ell_k K_k$ とすれば $a_i=a/\ell_i$、$b_i=b/\ell_i$ となる。
$i\ne j$ の対角比較で $[K_i,D_i]=H_i$ と $H_j$ が独立なので、
$X_i,Y_i$ の対称部分は各々 $h_iD_i,m_iD_i$ に限られる。
$[K_i,D_j]=-[K_j,D_i]\ne0$ より $a m_j\ell_j=-b h_i\ell_i$。
従って $h_i=h/\ell_i,m_i=m/\ell_i$、$am=-bh$ である。
$C_{ii}=(am-bh)H_i/\ell_i^2$ は二次元を張り $\sum C_{ii}=0$ なので、
三つの $\ell_i^2$ は等しい。符号の積が1の変更を対角共役で行い、三つの $\ell_i$ を等しくする。
残る二つの非零倍率の比を $(t^2-1)/t$ の非零定数倍で表せば、全体倍率を除き表示の $C(t)$ になる。

具体的に $C(t)_{ii}=-2t H_i$、$C(t)_{ij}=(1-t^2)[K_i,K_j]$ $(i\ne j)$。
$X_i=aK_i+hD_i$、$Y_i=bK_i+mD_i$ の全解は
\[
            bh=t,\quad am=-t,\quad ab+hm=1-t^2,\qquad (ab-1)(ab+t^2)=0 .
\]
両側の読み方に同じ計算が使える。元の $Q=K+tD,P=K-tD$ の選択を替えると
$Q=K-D/t,P=-t^2K-tD$ になり、元の $G=tK+D,F=tK-D$ の選択を替えると
逆倍率まで $G=K-tD,F=-(K+tD)$ になる。互いに独立な二倍率は残余群で吸収できる。
両方を替えた窓は $W(-1/t)$ である。実際 $C(-1/t)=-C(t)/t^2$ である。
$x=-1$ では各側の二根が一致するので集合として選択は一つである。

入力の形は枠の選択の存在で定義され、各実現上では枠自身が復元できる。
両読みの像五空間は対角跡零と交代行列の和で、その跡直交補は $D_i$ の包である。
この三空間の行列式は座標積の二倍であり、三角形の三頂点の行列の核が元の三座標軸になる。
従って枠を忘れても $i\ne j$ の
$\operatorname{tr}(C_{ij}C_{ji})/\operatorname{tr}(C_{ii}^2)=(x-1)^2/(16x)$ は同じである。
逆に $t\mapsto-t$ は二枠の同時奇置換、$t\mapsto-1/t$ は上記全体倍率で実現する。
\end{proof}

\begin{proof}[一方だけを替えた窓]
一つの側を替えた表示を $G=P=K-tD$、$Q=-F=K+tD$ とする。
先の一位置移動公式で $(G',F',Q',P')=(-P,-Q,-Q,P)$ となるから、
残余の $A=B=-I,s=-1$ で元に戻り最小周期1である。
最初の二位置を同じ四生成に読む際、第二位置の補三文字の符号を反転する。
するとその二次関係は、星文字 $h$ が中心で、巡回的な $(i,j,k)$ ごとに
\[
             (1-t)x_i x_j-(1+t)x_jx_i=-h x_k
\]
となる。実際 $Q_{b,ac}=-G_{a,bc}$ であり、二つの星を用いた六関係を足し引きすればよい。
他方だけの変更は $t\mapsto-1/t$ で得る。$a=1-t,b=1+t$ とし $ab\ne0$。
中心環 $k[h]$ 上で $yx,zy,zx$ を昇順に書き換えると、唯一の重なり $zyx$ は
両順とも $\frac ab xyz+\frac hb(x^2-y^2+z^2)$ になる。
補文字の長さ、次いで辞書式順序を用いれば停止する。従って $h$ は非零因子であり、
その商はパラメータ $q=b/a$ の三生成量子多項式環、中心一次拡大は標準Koszul左右AS正則である。
$q\ne1$。中心文字を規格化し生成元を一斉に倍率変更すれば
$X_i X_j-q X_jX_i=(1-q)h X_k$ と書ける。

$h=0$ の点対応は量子多項式環のもので、三座標直線を含み射影平面を張る。
一方 $h\ne0$ なら中心性により両点は同じであり、$h=1$ で $X_i X_j=X_k$ を解くと
原点と積1の四符号点だけである。従って無印の添字保持同型でも、
正次元の点軌跡の包から無限遠平面が決まり、残る五点の置換は全位置で同じである。
五点は射影枠で、原点は四点の重心という唯一のアフィン従属で識別される。
そのため全位置の射影変換は一致し、四符号点を保つ線型変換は積1の符号と三軸置換の全二十四個である。
巡回の向きを保てば $q'=q$、逆にすれば $q'=q^{-1}$ で、直接関係を送れば両者は十分でもある。
これは $s=\pm t$ と同値である。$W$ 同士の同型と正則性は
定理\ref{thm:marked-mixed-curve-isomorphisms}と命題\ref{prop:split-mixed-transverse-curve}による。
通常の $W$ と $A$ は最小周期で区別され、$x=-1$ での一致は上の二根の合流による。
\end{proof}

\begin{proof}[対角像だけが残る端点]
$t=\pm1$ では $[X_i,Y_j]=0$ $(i\ne j)$、$[X_i,Y_i]=-2t H_i\ne0$ である。
各組とも三独立（他側の各元と交換子を取る）で、全 $X_i,Y_i$ は全 $H_j$ に跡直交し対角が零である。
$X_i$ が非巡回ならスカラーと階数一 $w b^t$ の和であり、
$H_i(\ker b^t)\subset kw$ である。$H_i$ の階数が2なので $\ker b^t$ はその核 $ke_i$ を含み、
$w\in\operatorname{im} H_i$ となる。従って $i$ 番目の成分によりスカラーは零、
残る二行二列でも対角が零なので $X_i$ は他の二軸間の行列単位の倍である。$Y_i$ も同様。
もし $X_i$ が巡回なら $Y_j,Y_k$ は $X_i$ の跡零中心化環の全体を張る。
$Y_j$ が巡回ならこの可換環がその中心化環でもあり、$X_k$ は $Y_j$ と交換して
$Y_k$ とも交換することになり矛盾する。よって $Y_j,Y_k$ はどちらも非巡回である。
上で決めた異なる二辺の行列単位が可換なら行か列が共通となり、その包は階数1以下で
巡回な $X_i$ を含めない。結局全元が上記の行列単位の倍である。

対角交換子の値により $X_i,Y_i$ は逆向きの単位で倍率の積が決まり、
非対角交換子の消失により $X_i$ の三矢は有向三角形またはその全反転となる。
両側の六倍率は枠の対角変更で全て吸収できる。
実際二つの向きを $\epsilon,\delta\in\{1,-1\}$ とし $R$ を巡回する隣接二座標の差の行列とすれば、
倍率トーラス準同型の指数行列の行列式は
$\det(I-\epsilon\delta R^2)=(1+3\epsilon\delta)^2\ne0$ である。
両方向を同時に反転する操作は二枠の同じ奇置換と全体倍率で実現する。
よって元の向きの組合せは $W(1)$、片側反転の組合せは上記の中心表示の $t=1$ 型に尽きる。
後者では $h$ を零にすれば一方向の三積だけが零になる単項式商を得る。
長さ $n\ge1$ の正規語は $3\cdot2^{n-1}$ 個なので元の窓も標準成長を持たない。
$W(1)$ は既知の正則点である。その第一辺は有限なので他の $A(s)$ とも同型にならない。
実際点を $[r:x_1:x_2:x_3],[s:y_1:y_2:y_3]$ とすると
$r y_i+2x_j y_k=s x_i+2x_k y_j=0$ が巡回的に成り立つ。
$r=0$ では座標軸三点だけ、$r\ne0$ では星点または全 $x_i\ne0$ で
$r^3=-8x_1x_2x_3$、$x_1^4=x_2^4=x_3^4$ となり有限である。
\path{computations/verify_five_dimensional_mixed_brackets.py} は二つの選択の恒等式、
十次元の跡条件空間の小行列式、中心表示の係数と端点の倍率指数を検算する。
\end{proof}
